% GENERATED FILE. Edit canonical lessons, then run npm run generate:books.
% canonical-source-hash: fnv1a64:2b25df46516f1e08
% canonical-lessons: SA-C369-bhagyam, SA-C369-daivam, SA-C369-akasmat, SA-C369-upalabhate, SA-C369-sambhavana

\chapter{Luck and Fate}
\label{ch:sa-luck-and-fate}
\hlchaptermodality{369}

\begin{chapteropening}
\textbf{By the end of this chapter:} \emph{I can say what luck or fate might bring.}

Complete the last lesson of chapter 369: I can say what luck or fate might bring.
\end{chapteropening}

\section[\texorpdfstring{\sk{भाग्यम्} (bhāgyam)}{भाग्यम् (bhāgyam)}]{\sk{भाग्यम्} (bhāgyam) — luck; fortune}

\label{lesson:SA-C369-bhagyam}



\begin{quote}
Before the new one: say the Sanskrit for an inference, then the Sanskrit for so they say.
\end{quote}

\subsection*{You'll want to know: \sk{भाग्यम्}}
\textbf{\sk{भाग्यम्}} — \emph{bhāgyam} — \textquotedblleft{}luck; fortune\textquotedblright{}.

\begin{quote}
\sk{यदि} \sk{भाग्यम्} \sk{अस्ति}, \sk{तर्हि} \sk{सर्वं} \sk{सुखं} \sk{भविष्यति।}
\end{quote}

\begin{quote}
\emph{yadi bhāgyam asti, tarhi sarvaṃ sukhaṃ bhaviṣyati.}
\end{quote}

\begin{quote}
If we have luck, all will be well.
\end{quote}

A neuter noun.

\begin{cousinweb}
\textbf{\sk{भाग्यम्}} (\emph{bhāgyam}) comes from the root √bhaj, \textquotedblleft{}to share\textquotedblright{} (MW takes \emph{bhāgya} as \textquotedblleft{}to be shared\textquotedblright{}), and MW gives the neuter \emph{bhāgya} \textquotedblleft{}fate, destiny, fortune, good fortune\textquotedblright{} (MW).
\end{cousinweb}

\begin{grammarlens}[title={What you've built}]
One more word for saying what happens if.
\end{grammarlens}

\subsection*{Your turn}
\begin{itemize}
\raggedright
  \item \emph{Say it:} \emph{bhāgyam}
  \item \emph{Say it:} \emph{bhāgyam}, once more
  \item \emph{Say it:} \emph{yadi bhāgyam asti, tarhi sarvaṃ sukhaṃ bhaviṣyati.} — If we have luck, all will be well.
  \item \emph{Recall:} say the Sanskrit for an inference, then the Sanskrit for so they say, then say \emph{bhāgyam} again
\end{itemize}

\begin{tcolorbox}[breakable,colback=teal!4,colframe=teal!35!black,title={Before you move on}]
What does \sk{भाग्यम्} mean? (\textquotedblleft{}Luck; fortune\textquotedblright{}.) Say it once more.
\end{tcolorbox}
\section[\texorpdfstring{\sk{दैवम्} (daivam)}{दैवम् (daivam)}]{\sk{दैवम्} (daivam) — fate}

\label{lesson:SA-C369-daivam}



\begin{quote}
Before the new one: say the Sanskrit for so they say, then the Sanskrit for luck.
\end{quote}

\subsection*{You'll want to know: \sk{दैवम्}}
\textbf{\sk{दैवम्}} — \emph{daivam} — \textquotedblleft{}fate\textquotedblright{}.

\begin{quote}
\sk{यदि} \sk{दैवम्} \sk{इच्छति}, \sk{तर्हि} \sk{सर्वं} \sk{भविष्यति।}
\end{quote}

\begin{quote}
\emph{yadi daivam icchati, tarhi sarvaṃ bhaviṣyati.}
\end{quote}

\begin{quote}
If fate wills it, everything will happen.
\end{quote}

A neuter noun.

\begin{cousinweb}
\textbf{\sk{दैवम्}} (\emph{daivam}) is built on \emph{deva}, \textquotedblleft{}a god\textquotedblright{}; MW gives \emph{daiva} \textquotedblleft{}divine power or will, destiny, fate, chance\textquotedblright{} (MW).
\end{cousinweb}

\begin{grammarlens}[title={What you've built}]
One more word for saying what happens if.
\end{grammarlens}

\subsection*{Your turn}
\begin{itemize}
\raggedright
  \item \emph{Say it:} \emph{daivam}
  \item \emph{Say it:} \emph{daivam}, once more
  \item \emph{Say it:} \emph{yadi daivam icchati, tarhi sarvaṃ bhaviṣyati.} — If fate wills it, everything will happen.
  \item \emph{Recall:} say the Sanskrit for so they say, then the Sanskrit for luck, then say \emph{daivam} again
\end{itemize}

\begin{tcolorbox}[breakable,colback=teal!4,colframe=teal!35!black,title={Before you move on}]
What does \sk{दैवम्} mean? (\textquotedblleft{}Fate\textquotedblright{}.) Say it once more.
\end{tcolorbox}
\section[\texorpdfstring{\sk{अकस्मात्} (akasmāt)}{अकस्मात् (akasmāt)}]{\sk{अकस्मात्} (akasmāt) — by accident; without warning}

\label{lesson:SA-C369-akasmat}



\begin{quote}
Before the new one: say the Sanskrit for luck, then the Sanskrit for fate.
\end{quote}

\subsection*{You'll want to know: \sk{अकस्मात्}}
\textbf{\sk{अकस्मात्}} — \emph{akasmāt} — \textquotedblleft{}by accident; without warning\textquotedblright{}.

\begin{quote}
\sk{यदि} \sk{अकस्मात्} \sk{वर्षा} \sk{भवति}, \sk{तर्हि} \sk{अहं} \sk{गृहे} \sk{स्थास्यामि।}
\end{quote}

\begin{quote}
\emph{yadi akasmāt varṣā bhavati, tarhi ahaṃ gṛhe sthāsyāmi.}
\end{quote}

\begin{quote}
If it rains without warning, I will stay at home.
\end{quote}

\begin{cousinweb}
\textbf{\sk{अकस्मात्}} (\emph{akasmāt}) is built from \emph{a-}, \textquotedblleft{}not\textquotedblright{}, and \emph{kasmāt}, \textquotedblleft{}from what? why?\textquotedblright{}; MW gives \emph{a-kasmāt} \textquotedblleft{}without a why or a wherefore, accidentally, suddenly\textquotedblright{} (MW).
\end{cousinweb}

\begin{grammarlens}[title={What you've built}]
One more word for saying what happens if.
\end{grammarlens}

\subsection*{Your turn}
\begin{itemize}
\raggedright
  \item \emph{Say it:} \emph{akasmāt}
  \item \emph{Say it:} \emph{akasmāt}, once more
  \item \emph{Say it:} \emph{yadi akasmāt varṣā bhavati, tarhi ahaṃ gṛhe sthāsyāmi.} — If it rains without warning, I will stay at home.
  \item \emph{Recall:} say the Sanskrit for luck, then the Sanskrit for fate, then say \emph{akasmāt} again
\end{itemize}

\begin{tcolorbox}[breakable,colback=teal!4,colframe=teal!35!black,title={Before you move on}]
What does \sk{अकस्मात्} mean? (\textquotedblleft{}By accident; without warning\textquotedblright{}.) Say it once more.
\end{tcolorbox}
\section[\texorpdfstring{\sk{उपलभते} (upalabhate)}{उपलभते (upalabhate)}]{\sk{उपलभते} (upalabhate) — comes by; obtains}

\label{lesson:SA-C369-upalabhate}



\begin{quote}
Before the new one: say the Sanskrit for fate, then the Sanskrit for by accident.
\end{quote}

\subsection*{You'll want to know: \sk{उपलभते}}
\textbf{\sk{उपलभते}} — \emph{upalabhate} — \textquotedblleft{}comes by; obtains\textquotedblright{}.

\begin{quote}
\sk{यदि} \sk{भाग्यम्} \sk{अस्ति}, \sk{तर्हि} \sk{सः} \sk{धनम्} \sk{उपलभते।}
\end{quote}

\begin{quote}
\emph{yadi bhāgyam asti, tarhi saḥ dhanam upalabhate.}
\end{quote}

\begin{quote}
If he is lucky, he comes by money.
\end{quote}

Past (he …): \sk{सः} \sk{धनम्} \sk{उपलब्धवान्}, \emph{saḥ dhanam upalabdhavān}.

\begin{cousinweb}
\textbf{\sk{उपलभते}} (\emph{upalabhate}) is built on \emph{upa}, \textquotedblleft{}towards\textquotedblright{}, and the root √labh, \textquotedblleft{}to take, to get\textquotedblright{}; MW gives \emph{upa-√labh} \textquotedblleft{}to acquire, to receive, to obtain, to find\textquotedblright{} (MW).
\end{cousinweb}

\begin{grammarlens}[title={What you've built}]
One more word for saying what happens if.
\end{grammarlens}

\subsection*{Your turn}
\begin{itemize}
\raggedright
  \item \emph{Say it:} \emph{upalabhate}
  \item \emph{Say it:} \emph{upalabhate}, once more
  \item \emph{Say it:} \emph{yadi bhāgyam asti, tarhi saḥ dhanam upalabhate.} — If he is lucky, he comes by money.
  \item \emph{Recall:} say the Sanskrit for fate, then the Sanskrit for by accident, then say \emph{upalabhate} again
\end{itemize}

\begin{tcolorbox}[breakable,colback=teal!4,colframe=teal!35!black,title={Before you move on}]
What does \sk{उपलभते} mean? (\textquotedblleft{}Comes by; obtains\textquotedblright{}.) Say it once more.
\end{tcolorbox}
\section[\texorpdfstring{\sk{सम्भावना} (sambhāvanā)}{सम्भावना (sambhāvanā)}]{\sk{सम्भावना} (sambhāvanā) — a likelihood}

\label{lesson:SA-C369-sambhavana}



\begin{quote}
Before the new one: say the Sanskrit for by accident, then the Sanskrit for comes by.
\end{quote}

\subsection*{You'll want to know: \sk{सम्भावना}}
\textbf{\sk{सम्भावना}} — \emph{sambhāvanā} — \textquotedblleft{}a likelihood\textquotedblright{}.

\begin{quote}
\sk{यदि} \sk{आकाशः} \sk{कृष्णः}, \sk{तर्हि} \sk{वर्षायाः} \sk{सम्भावना} \sk{अस्ति।}
\end{quote}

\begin{quote}
\emph{yadi ākāśaḥ kṛṣṇaḥ, tarhi varṣāyāḥ sambhāvanā asti.}
\end{quote}

\begin{quote}
If the sky is dark, rain is likely.
\end{quote}

A feminine noun.

\begin{cousinweb}
\textbf{\sk{सम्भावना}} (\emph{sambhāvanā}) is built on \emph{sam} and the root √bhū, \textquotedblleft{}to be\textquotedblright{}, in its causative form; MW gives \emph{sam-bhāvanā} \textquotedblleft{}imagination, supposition, assumption\textquotedblright{}, and \textquotedblleft{}a likelihood\textquotedblright{} is its modern sense (MW).
\end{cousinweb}

\begin{grammarlens}[title={What you've built}]
One more word for saying what happens if.
\end{grammarlens}

\subsection*{Your turn}
\begin{itemize}
\raggedright
  \item \emph{Say it:} \emph{sambhāvanā}
  \item \emph{Say it:} \emph{sambhāvanā}, once more
  \item \emph{Say it:} \emph{yadi ākāśaḥ kṛṣṇaḥ, tarhi varṣāyāḥ sambhāvanā asti.} — If the sky is dark, rain is likely.
  \item \emph{Recall:} say the Sanskrit for by accident, then the Sanskrit for comes by, then say \emph{sambhāvanā} again
\end{itemize}

\begin{tcolorbox}[breakable,colback=teal!4,colframe=teal!35!black,title={Before you move on}]
What does \sk{सम्भावना} mean? (\textquotedblleft{}A likelihood\textquotedblright{}.) Say it once more.
\end{tcolorbox}
